% !TEX program = xelatex
% Fixture: requiere XeLaTeX (fontspec falla con pdfLaTeX).
% Las fuentes se cargan por nombre de archivo para no depender de fontconfig.
\documentclass[11pt,a4paper]{article}
\usepackage{fontspec}
\usepackage{unicode-math}
\usepackage[spanish]{babel}

\setmainfont{texgyrepagella}[
  Extension      = .otf,
  UprightFont    = *-regular,
  BoldFont       = *-bold,
  ItalicFont     = *-italic,
  BoldItalicFont = *-bolditalic,
]
\setsansfont{texgyreheros}[
  Extension      = .otf,
  UprightFont    = *-regular,
  BoldFont       = *-bold,
  ItalicFont     = *-italic,
  BoldItalicFont = *-bolditalic,
]
\setmonofont{lmmono10-regular.otf}
\setmathfont{latinmodern-math.otf}

\title{Fuentes OpenType con XeLaTeX}
\author{Fixture de pruebas}
\date{}

\begin{document}
\maketitle

Texto en TeX Gyre Pagella con \textbf{negrita}, \textit{cursiva} y
\textbf{\textit{ambas}}. Caracteres Unicode escritos directamente: ñ, ü, €, —,
«comillas», “tipográficas” y ligaduras: fi, fl, ffi.

\textsf{Texto sin serifa en TeX Gyre Heros.}

\texttt{Texto monoespaciado en Latin Modern Mono.}

Matemática Unicode: $α + β = γ$, $∑_{i=1}^{n} i = \frac{n(n+1)}{2}$, $x ≤ y$.

\end{document}
